\chapter{\IfLanguageName{dutch}{Stand van zaken}{State of the art}}
\label{ch:stand-van-zaken}

In dit hoofdstuk wordt een overzicht gegeven van de huidige stand van zaken binnen het domein van journalistieke partijdigheid en automatische detectie van bias. De literatuurstudie bouwt voort op de probleemstelling en onderzoeksvragen beschreven in Hoofdstuk \ref{ch:inleiding} en dient als basis voor de keuzes die in Hoofdstuk \ref{ch:methodologie} worden besproken.

\section{Vormen van journalistieke vooringenomenheid}
\label{sec:vormen-van-journalistieke-vooringenomenheid}

Partijdigheid komt voor in verschillende vormen, elk met specifieke kenmerken en effecten op de inhoud van nieuwsartikelen. Deze vormen van partijdigheid kunnen zichtbaar worden in woordkeuze, selectie van bronnen, nadruk op bepaalde standpunten of het weglaten van relevante informatie. Inzicht in de verschillende vormen is essentieel voor de ontwikkeling van automatische detectiesystemen.

In de literatuur wordt journalistieke bias in verschillende vormen geclassificeerd. \textcite{Mullainathan2002} onderscheiden twee categorieën: \textit{ideology bias} en \textit{spin bias}, terwijl \textcite{SaezTrumper2013} drie andere categorieën onderscheiden: \textit{gatekeeping}, \textit{coverage bias} en \textit{statement bias}.

\subsection{Ideology bias}

Bij \textit{ideology bias} wordt de inhoud van een artikel gestuurd in een bepaalde ideologische richting. Vaak reflecteren deze ideologische richtingen de voorkeuren van het mediabedrijf en haar doelgroep \autocite{Mullainathan2002}. Het sturen van informatie in bepaalde politieke richtingen is ook waargenomen bij VTM Nieuws, zoals gerapporteerd in het onpartijdigheidsrapport van de \gls{vrm} \autocite{VanAelst2024}.

Recente studies tonen aan dat de relatie tussen de ideologische voorkeur van journalisten en hun berichtgeving niet altijd duidelijk is. Hoewel journalisten gemiddeld vaker een linksere visie hebben, blijft de inhoud van nieuwsartikelen vaak neutraal. Dit komt doordat journalisten door objectiviteit hun eigen mening proberen te beperken \autocite{Coddington2026}. Ideologische voorkeuren komen daardoor vaker naar voren in subtielere vormen, zoals woordkeuze en framing, dan in een duidelijke politieke richting.

\subsection{Spin bias}

In het geval van \textit{spin bias} wordt informatie gemanipuleerd met als doel het artikel opvallender of sensationeler te maken. Door de concurrentie tussen nieuwsorganisaties worden journalisten aangemoedigd om opvallende artikelen te schrijven, wat in sommige gevallen kan leiden tot het verdraaien of overdrijven van feiten \autocite{Alsem2008}. Er is nog geen onderzoek uitgevoerd naar het aandeel van \textit{spin bias} in artikelen gepubliceerd door \gls{hln}.

\subsection{Gatekeeping}

\textit{Gatekeeping}, ook wel \textit{selection bias} genoemd, wordt gedefinieerd als het bepalen van wie bepaalde informatie, macht of kansen krijgt \autocite{CALD2025}. De auteur maakt een specifieke selectie van informatie voor het artikel en laat andere informatie achterwege. Voor nieuwsartikelen worden auteurs vaak gedwongen om een selectie van informatie te maken vanwege de beperkte ruimte \autocite{SaezTrumper2013}.

\subsection{Coverage bias}

Bij \textit{coverage bias} wordt meer aandacht toegekend aan bepaalde onderdelen of personen in het artikel, en worden andere onderdelen of personen onderdrukt \autocite{DAlessio2006}. Deze extra aandacht kan zich op verschillende manieren uiten, zoals meer zendtijd op televisie, meer spreektijd op de radio of meer tekst in artikelen. Meer spreektijd geven aan bepaalde politieke personen is ook waargenomen bij VTM Nieuws \autocite{VanAelst2024}.

\subsection{Statement bias}

\textit{Statement bias}, ook wel \textit{presentation bias} genoemd, verwijst naar het voordeliger omschrijven van bepaalde onderdelen of personen, ten koste van andere \autocite{SaezTrumper2013}. Volgens \textcite{DAlessio2006} kan een mediabericht pas beschouwd worden als neutraal wanneer het evenveel uitspraken bevat die in de ene richting voordelig zijn als uitspraken die in de tegenovergestelde richting voordelig zijn.

\section{Methoden om bias te detecteren}
\label{sec:methoden-om-bias-te-detecteren}

Het identificeren van journalistieke bias in nieuwsartikelen vereist methoden die zowel tekstuele als semantische kenmerken kunnen detecteren. Het analyseren van deze kenmerken vormt de basis voor de uiteindelijke classificatie van artikelen, waarbij technieken uit het domein van \gls{nlp} kunnen worden toegepast.

\subsection{Supervised learning}

Een van de methoden om artikelen te classificeren op basis van partijdigheid is \textit{supervised learning}, waarbij modellen getraind worden op gelabelde data. Vaak worden neurale tekstclassificatoren zoals \gls{bert} hierbij gebruikt.

\gls{bert} is een complexe architectuur die contextuele representaties van teksten kan genereren en heeft zich al bewezen als een succesvolle methode voor \gls{nlp}-taken \autocite{Baly2020}. \gls{bert} is gebaseerd op het encoder-gedeelte van de transformer-architectuur \autocite{Devlin2018}, zoals weergegeven in figuur \ref{fig:bert-architectuur} \autocite{Vaswani2017}.

In het onderzoek van \textcite{Baly2020} werd een biasclassificatiemodel op basis van \gls{bert} ontwikkeld. Het model werd getest op artikelen uit kranten die niet in de trainingsset waren opgenomen, om te voorkomen dat de nieuwsbron herkend zou worden en de classificatie zou beïnvloeden. In dit onderzoek zal hier geen rekening mee gehouden worden, aangezien het beperkt is tot één enkele nieuwsbron.

Er bestaan verschillende voorgetrainde \gls{bert}-modellen, beschikbaar via \textit{Hugging Face} \autocite{Valurank2022}. \gls{bert} kan eenvoudig worden aangeroepen als een black-box model, waarbij de invoer een tekst is en de uitvoer een classificatie van de partijdigheid van die tekst. Een voorbeeld van een minimale implementatie van \gls{bert} voor tekstclassificatie in Python is te zien in codefragment \ref{code:bert-python-voorbeeld}.

\begin{figure}
  \centering
  \includegraphics[width=0.2\textwidth]{bert-architectuur.png}
  \caption[\gls{bert}-archi\-tectuur]{\label{fig:bert-architectuur}Schematische weergave van het encoder-gedeelte van de transformer-architectuur, waarop \gls{bert} is gebaseerd. Bron: \textcite{Vaswani2017}}
\end{figure}

\begin{listing}
  \begin{minted}{python}
    from transformers import pipeline

    classifier = pipeline("text-classification", model="valurank/finetuned-distilbert-news-article-categorization")

    text = "The irony, of course, is that the exhibit that invites people to throw trash at vacuuming Ivanka Trump lookalike reflects every stereotype feminists claim to stand against, oversexualizing Ivanka's body and ignoring her hard work."

    prediction = classifier(text)
  \end{minted}
  \caption[Voorbeeld \gls{bert}-inference]{Voorbeeld van het aanroepen van een voorgetraind \gls{bert}-model voor tekstclassificatie. Het model voorspelt de partijdigheid van een korte Engelse tekst. Bron: \textcite{Valurank2022}.}
  \label{code:bert-python-voorbeeld}
\end{listing}

\subsection{Unsupervised learning}

Een andere methode die gebruikt kan worden is \textit{unsupervised learning}, waarbij de dataset niet gelabeld is. Hierbij wordt vaak gebruik gemaakt van vectorisatiemethoden in combinatie met clusteringalgoritmen \autocite{UmarNadeem2022}. Uit een onderzoek van \textcite{UmarNadeem2022} werd geconcludeerd dat deze methode niet effectief is voor het vinden van tekstpatronen die essentieel zijn voor het detecteren van bias. Figuur \ref{fig:unsupervised-learning-scheme} illustreert het verschil tussen supervised en unsupervised learning.

\begin{figure}
  \centering
  \includegraphics[width=0.6\textwidth]{unsupervised_vs_supervised.png}
  \caption[Supervised vs Unsupervised learning]{Schematische weergave van het verschil in werking tussen supervised en unsupervised learning. Bron: \textcite{Morimoto2021}.}
  \label{fig:unsupervised-learning-scheme}
\end{figure}

\subsection{Sentiment analysis}

Een derde methode die relevant zou kunnen zijn binnen ons onderzoek is \textit{sentimentanalyse}. Deze techniek richt zich op het identificeren van toon in teksten en kan beoordelen of deze positief, negatief of neutraal is \autocite{Cui2025}. Hoewel deze methode mogelijk niet voldoende is om alle vormen van bias te detecteren, kan het wel in combinatie met een \textit{supervised learning} techniek zoals \gls{bert} gebruikt worden om betere resultaten te bereiken. Een voorbeeld van een eenvoudige implementatie van sentimentanalyse in Python is te zien in codefragment \ref{code:sentimentanalyse-voorbeeld}.

\begin{listing}
  \begin{minted}{python}
    from transformers import pipeline

    classifier = pipeline("sentiment-analysis")

    text = "The irony, of course, is that the exhibit that invites people to throw trash at vacuuming Ivanka Trump lookalike reflects every stereotype feminists claim to stand against, oversexualizing Ivanka's body and ignoring her hard work."

    prediction = classifier(text)
  \end{minted}
  \caption[Voorbeeld sentimentanalyse]{Voorbeeld van het uitvoeren van sentimentanalyse op een korte Engelse tekst. Het model classificeert de toon van de tekst als positief of negatief.}
  \label{code:sentimentanalyse-voorbeeld}
\end{listing}

\subsection{Rule-based approach}

\textit{Rule-based benaderingen} maken gebruik van vooraf gedefinieerde regels om bias-gerelateerde kenmerken in teksten te identificeren. In een recente studie van \textcite{Martinez2024} werd een rule-based model vergeleken met een deep learning model voor politieke biasclassificatie. Hieruit bleek dat het rule-based model consistent presteert en een hogere transparantie biedt. Deze studie toont aan dat rule-based benaderingen waardevol kunnen blijven naast complexe neurale modellen.

\section{Methodes om tekst te extraheren uit een webpagina}
\label{sec:methodes-om-tekst-te-extraheren-uit-een-webpagina}

Voor de ontwikkeling van de browserextensie is het noodzakelijk om de tekst van geopende nieuwsartikelen op de \gls{hln}-website te kunnen extraheren. Een veelgebruikte methode hiervoor is het doorlopen van de \gls{dom}-structuur van de webpagina \autocite{GCD2012}. Het \gls{dom} vormt een hiërarchische representatie van alle elementen op een webpagina, zie figuur \ref{fig:dom-structuur}, en maakt het mogelijk om deze elementen dynamisch te lezen en te wijzigen. Moderne websites gedragen zich doorgaans consistent binnen verschillende browsers, waardoor het element dat de hoofdinhoud bevat meestal kan worden geïdentificeerd als het element met het grootste tekstvolume op een vergelijkbare diepte in de \gls{dom}-structuur \autocite{Aleksandersen2018}.

Hoewel er alternatieve methodes bestaan om tekstextractie te bereiken, wordt dit onderzoek beperkt tot het gebruik van het \gls{dom}. De keuze wordt gemotiveerd door het feit dat dit een goed ondersteunde methode is binnen de meest gebruikte browsers.

\begin{figure}
  \centering
  \includegraphics[width=0.8\textwidth]{dom-tree.png}
  \caption[Schematische \gls{dom}-structuur]{Schematische voorstelling van een vereenvoudigde \gls{dom}-structuur. Bron: \textcite{W3Schools2024}.}
  \label{fig:dom-structuur}
\end{figure}

\section{Methodes om semantische betekenis te behouden}
\label{sec:methodes-om-semantische-betekenis-te-behouden}

Het behoud van semantische betekenis van woorden is cruciaal binnen dit onderzoek. Een veelgebruikte methode hiervoor is het gebruik van woordembeddings, waarbij woorden worden omgezet in vectorrepresentaties. Volgens \textcite{Akbik2018} kan binnen deze benadering onderscheid worden gemaakt tussen klassieke en contextuele woordembeddings.

\subsection{Global Vector Representations}

Klassieke woordembeddings koppelen ieder woord aan één vaste vector, onafhankelijk van de context waarin het voorkomt. Een bekende techniek binnen deze categorie is \gls{glove}, dat globale corpusstatistieken gebruikt in combinatie met \textit{unsupervised learning} om woorden in een vectorruimte te representeren \autocite{Pennington2014}. Deze modellen kunnen semantische relaties vastleggen, maar houden geen rekening met contextuele betekenissen van individuele woorden.

\subsection{Contextualized Word Representations}

Om ook contextuele informatie te behouden binnen een vectorrepresentatie, kunnen contextuele woordembeddings worden gebruikt. Deze embeddings genereren representaties die afhankelijk zijn van de context waarin een woord voorkomt binnen een zin \autocite{Akbik2018}. Een belangrijke techniek binnen deze categorie is het gebruik van transformer-gebaseerde representaties zoals \gls{bert} \autocite{Khodake2020}, die contextuele informatie kunnen vasthouden in een vectorrepresentatie.

\subsection{Embeddings from Language Models}

Een andere relevante techniek is het gebruik van \gls{elmo}, waarbij vectorrepresentaties van woorden worden afgeleid uit een bidirectioneel \gls{lstm}-taalmodel \autocite{Peters2018}. \gls{elmo} combineert informatie uit alle interne lagen van het taalmodel tot een gewogen vectorrepresentatie voor elk woord, waardoor zowel semantische als contextuele betekenissen van woorden kunnen worden vastgelegd.

\section{Samenstellen van een dataset}
\label{sec:samenstellen-van-een-dataset}

Het is niet mogelijk om een grootschalige, gelabelde dataset van nieuwsartikelen van \gls{hln} te verkrijgen. Handmatige labeling is bovendien niet haalbaar, vanwege de complexiteit en tijdsintensiviteit van het beoordelen van bias in nieuwsartikelen. Daarom is er nood aan een trainingsstrategie die deze beperkingen kan omzeilen.

\begin{enumerate}
    \item \textit{Metadata-dataset}: Een metadata-dataset kan verzameld worden via webscraping van de \gls{hln}-website en bevat aanvullende informatie over nieuwsartikelen, zoals publiek beschikbare gegevens over de auteur(s). Deze metadata kan gebruikt worden om patronen in partijdigheid over meerdere publicaties van een auteur te analyseren en kan als aanvullende context dienen binnen een biasclassificatiemodel.
    \item \textit{Sentence-level dataset}: Op zinsniveau gelabelde datasets, zoals de \textit{FactNews}-dataset \autocite{Vargas2024}, voorgesteld in het onderzoek van \textcite{Vargas2023}, kunnen gebruikt worden om modellen te trainen die indicaties van bias binnen individuele zinnen kan detecteren. Deze datasets maken het mogelijk om zinnen te identificeren die sterk bijdragen aan de uiteindelijke biasclassificatie van een artikel.
    \item \textit{Document-level dataset}: Voor classificatie op artikelniveau kan gebruik worden gemaakt van bestaande datasets zoals de \textit{NewsMediaBias-Plus}-dataset \autocite{Raza2024}. Deze datasets kunnen gebruikt worden voor het trainen van modellen die globale bias van artikelen detecteren.
\end{enumerate}

Het is mogelijk om deze verschillende datasets te combineren binnen een pipeline, waarbij een globale biasclassificatie van een artikel wordt aangevuld met aanduidingen van zinnen die het meest bijdragen aan deze classificatie.

\section{Preprocessing van data}
\label{sec:preprocessing-van-data}

Vooraleer een machine learning-model kan worden getraind, is het noodzakelijk om de data te preprocessen. Hierbij wordt de data getransformeerd naar een formaat waaruit het model effectief kan leren. Aangezien dit onderzoek \gls{nlp}-taken omvat, is de preprocessing van tekstuele data een cruciale stap.

\begin{itemize}
    \item \textit{Tokenization} is het proces waarbij een tekst wordt opgesplitst in kleinere eenheden, zoals woorden of zinnen, die tokens worden genoemd. Hierdoor kunnen individuele woorden of zinnen worden geanalyseerd en verwerkt door het model.
    \item \textit{Stop-word removal} is het proces waarbij veelvoorkomende woorden die weinig betekenis toevoegen, zoals lidwoorden en voorzetsels, worden verwijderd. Dit helpt om de focus te leggen op de betekenisvollere woorden.
    \item \textit{Lowercase conversion} is het proces waarbij alle letters in een tekst worden omgezet naar kleine letters. Hierdoor hoeft het model geen onderscheid te maken tussen woorden zoals ``Nieuws'' en ``nieuws''.
    \item \textit{Stemming} is het proces waarbij woorden worden teruggebracht tot hun stam. De woorden ``loopt'' en ``lopen'' worden bijvoorbeeld teruggebracht naar ``loop''. Dit vermindert variatie in woordvormen en helpt het model om gelijkaardige woorden als één woord te behandelen \autocite{Uysal2014}.
    \item \textit{Lemmatization} is een gelijkaardig proces waarbij woorden worden herleid tot hun correcte woordenboekvorm. Het woord ``was'' wordt bijvoorbeeld herleid naar ``zijn''. In tegenstelling tot stemming houdt lemmatization rekening met de betekenis en grammaticale context van de woorden \autocite{Islam2025}.
\end{itemize}

\section{Splitsen van data}
\label{sec:splitsen-van-data}

Het opsplitsen van een dataset in afzonderlijke subsets is een belangrijke stap binnen het machine learning proces. Een correcte datasplit maakt het mogelijk om modellen te trainen, te valideren en te evalueren.

\subsection{Train-, validatie- en testset}

Bij een train-validatie-test splitsing wordt de dataset verdeeld in drie subsets:

\begin{itemize}
    \item \textit{Trainingsset}: wordt gebruikt om het model te trainen.
    \item \textit{Validatieset}: wordt gebruikt om de prestaties van het model te monitoren tijdens het trainingsproces.
    \item \textit{Testset}: wordt gebruikt voor de finale evaluatie van het model.
\end{itemize}

\subsection{Stratified sampling}

Bij classificatieproblemen met een ongelijke klassenverdeling kan \textit{stratified sampling} worden toegepast. Hierbij wordt ervoor gezorgd dat de verhouding tussen de verschillende klassen behouden blijft in elke subset van de dataset. Dit voorkomt dat bepaalde klassen ondervertegenwoordigd zijn in de train-, validatie- of testset.

\subsection{Cross-validation}

Bij \textit{k-fold cross-validation} wordt de dataset opgesplitst in $k$ gelijke delen. Het model wordt nadien meerdere keren getraind, waarbij telkens één deel wordt gebruikt als validatieset en de overige delen als trainingsset. Op deze manier wordt elk deel van de dataset één keer gebruikt voor validatie. Dit zorgt voor een betrouwbaarder beeld van de prestaties van het model.

\section{Evaluatiemethoden}
\label{sec:evaluatiemethoden}

Het evalueren van modellen voor media-biasclassificatie kan op verschillende manieren worden uitgevoerd. Het is essentieel om evaluatiemethoden te selecteren die aansluiten bij de doelstellingen van dit onderzoek.

\subsection{Klassieke classificatiemetrieken}

Voor de evaluatie van classificatiemodellen kunnen standaardmetrieken uit de machine learning-literatuur worden toegepast, waaronder \textit{accuracy}, \textit{precision}, \textit{recall} en de \textit{F1-score}. Deze metrieken bieden elk een ander perspectief op de prestaties van het model.

\begin{itemize}
    \item \textit{Accuracy} meet het percentage correct geclassificeerde voorbeelden. Deze metriek is eenvoudig te interpreteren, maar kan misleidend zijn bij een ongelijke klassenverdeling. Hierdoor kan de score hoog zijn wanneer het model goed presteert op de dominante klassen, terwijl het minder goed presteert op de minder frequente klassen.
    \item \textit{Precision} meet de hoeveelheid correct voorspelde positieve gevallen. Deze metriek is relevant in scenario's waar het beperken van \textit{false positives} belangrijk is.
    \item \textit{Recall} meet de hoeveelheid werkelijke positieve gevallen dat correct wordt voorspeld. Deze metriek is vooral belangrijk wanneer het vermijden van \textit{false negatives} cruciaal is.
    \item De \textit{F1-score} is het harmonisch gemiddelde van precision en recall en biedt een gebalanceerde maat voor de prestaties van het model. Deze metriek is nuttig wanneer een trade-off tussen \textit{precision} en \textit{recall} moet worden gemaakt.
    \item De \textit{macro F1-score} berekent het gemiddelde van de F1-scores over alle klassen, zonder rekening te houden met de frequentie van elke klasse. Hierdoor weegt elke klasse even zwaar mee, wat nuttig is bij datasets met een ongelijke klassenverdelingen.
\end{itemize}

\subsection{MBIB benchmark}

Naast klassieke evaluatiemetrieken kan gebruik worden gemaakt van benchmarks die specifiek zijn ontworpen voor media-biasdetectie. \gls{mbib} is een benchmark die meerdere datasets combineert om de prestaties van media-biasclassificatiemodellen te evalueren \autocite{Wessel2023}.

Deze benchmark bestaat uit datasets die verschillende vormen van bias bevatten, waaronder linguïstische, politieke en contextuele vormen van bias. Hierdoor maakt \gls{mbib} het mogelijk om de generalisatie van modellen te testen, waarbij prestaties worden gemeten aan de hand van de F1-score op zowel individuele datasets als een gecombineerde score over alle datasets.

\section{Overfitting en underfitting}
\label{sec:overfitting-en-underfitting}

Bij het trainen van machine learning-modellen is het essentieel om een evenwicht te vinden tussen het leren van voldoende relevante patronen en het vermijden dat het model de trainingsdata te letterlijk onthoudt. Twee veelvoorkomende problemen zijn \textit{overfitting} en \textit{underfitting}.

\begin{itemize}
    \item \textit{Overfitting} ontstaat wanneer een model te sterk wordt afgestemd op de trainingsdata, waardoor het niet goed generaliseert naar nieuwe data. Er zijn verschillende methoden om overfitting te beperken:
    \begin{itemize}
        \item \textit{Regularisatie} is een methode die complexe patronen bestraft en het model stimuleert om algemenere patronen te leren.
        \item \textit{Dropout} is een techniek waarbij tijdens het trainen willekeurig bepaalde onderdelen van het model tijdelijk worden uitgeschakeld, zodat het model niet afhankelijk wordt van specifieke onderdelen van de data.
        \item \textit{Early stopping} is een methode waarbij het trainen wordt gestopt zodra de prestaties op een aparte validatieset verslechteren.
        \item \textit{Cross-validation} is een techniek waarbij de dataset in meerdere subsets wordt verdeeld en het model meerdere keren wordt getraind en geëvalueerd op verschillende combinaties.
    \end{itemize}
    \item \textit{Underfitting} ontstaat wanneer een model onvoldoende complex is om relevante patronen te leren, wat leidt tot lage prestaties op zowel trainingsdata als nieuwe data. Mogelijke oplossingen zijn het gebruik van een complexer model en het toevoegen van meer of relevantere features.
\end{itemize}

\section{Bestaande libraries en tools}
\label{sec:bestaande-libraries-en-tools}

Voor de implementatie en evaluatie van media-bias classificatiemodellen kan gebruik worden gemaakt van verschillende bestaande libraries en tools.

\begin{itemize}
    \item \textit{Scikit-learn} is een veelgebruikte Python-library voor machine learning en biedt een uitgebreide set aan algoritmen en evaluatietools. De library bevat implementaties voor classificatie en preprocessing, en voorziet functionaliteiten voor evaluatie van modellen, zoals het berekenen van metrieken, het splitsen van datasets en het uitvoeren van cross-validation.
    \item \textit{Hugging Face Transformers} is een library die toegang biedt tot een grote collectie van voorgetrainde transformer-modellen die ingezet kunnen worden voor verschillende \gls{nlp}-taken. Deze modellen maken gebruik van deep learning-architecturen en genereren contextuele representaties van tekst, wat essentieel is voor het detecteren van bias.
    \item \textit{Pandas} is een Python-library voor data-analyse en -manipulatie en wordt vaak gebruikt voor het verwerken van datasets. \textit{Pandas} maakt het mogelijk om data efficiënt te transformeren en analyseren, wat handig is tijdens het preprocessen van data en het analyseren van modelresultaten.
    \item \textit{Matplotlib} en \textit{Seaborn} zijn Python-libraries voor datavisualisatie. Ze bieden functionaliteiten voor het genereren van grafieken en visualisaties, wat nuttig is voor het analyseren van modelprestaties en het interpreteren van resultaten.
    \item \textit{BeautifulSoup} is een Python-library voor het parsen van HTML-documenten. Deze library wordt gebruikt voor het extraheren van tekst uit webpagina's, wat essentieel is voor het verzamelen van nieuwsartikelen.
\end{itemize}